% =============================================================
%  Fossil Manifolds and Mediated Extraction:
%  LAP1-B, Regulated Instability, and the Geometry of Observability
%
%  Author: Flyxion
%  Engine: LuaLaTeX
% =============================================================

\documentclass[12pt, a4paper]{article}

% ── Font and typography ──────────────────────────────────────
\usepackage{fontspec}
\setmainfont{Linux Libertine O}[
  Ligatures = TeX,
  Numbers   = OldStyle,
]
\setsansfont{Linux Biolinum O}[Ligatures = TeX]
\setmonofont{DejaVu Sans Mono}[Scale = 0.88]

% ── Mathematics ─────────────────────────────────────────────
\usepackage{amsmath, amssymb, amsthm}
\usepackage{mathtools}

% ── Microtype and spacing ────────────────────────────────────
\usepackage{microtype}
\usepackage{setspace}
\setstretch{1.18}

% ── Page geometry ────────────────────────────────────────────
\usepackage[
  a4paper,
  top    = 3.2cm,
  bottom = 3.5cm,
  left   = 3.4cm,
  right  = 3.4cm,
  marginparwidth = 2.0cm,
  headheight = 13.6pt,
]{geometry}

% ── Section formatting ───────────────────────────────────────
\usepackage{titlesec}
\titleformat{\section}
  {\normalfont\large\bfseries\sffamily}
  {\thesection}{1em}{}
\titleformat{\subsection}
  {\normalfont\normalsize\bfseries\sffamily}
  {\thesubsection}{1em}{}
\titlespacing*{\section}{0pt}{2.4ex plus .6ex minus .2ex}{1.0ex plus .2ex}
\titlespacing*{\subsection}{0pt}{1.8ex plus .4ex minus .2ex}{0.8ex plus .2ex}

% ── Headers and footers ─────────────────────────────────────
\usepackage{fancyhdr}
\pagestyle{fancy}
\fancyhf{}
\fancyhead[L]{\small\sffamily Flyxion}
\fancyhead[R]{\small\sffamily Fossil Manifolds and Mediated Extraction}
\fancyfoot[C]{\small\thepage}
\renewcommand{\headrulewidth}{0.35pt}

% ── Hyperlinks ───────────────────────────────────────────────
\usepackage[
  colorlinks = true,
  linkcolor  = black,
  citecolor  = black,
  urlcolor   = black,
]{hyperref}

% ── References ───────────────────────────────────────────────
% Bibliography is embedded as \thebibliography at end of document.
% Use \citep{key} for parenthetical and \citet{key} for textual citations.
\usepackage{natbib}

% ── Epigraph ─────────────────────────────────────────────────
\usepackage{epigraph}
\setlength{\epigraphwidth}{0.72\textwidth}
\renewcommand{\epigraphflush}{center}
\renewcommand{\sourceflush}{center}

% ── Custom commands ──────────────────────────────────────────
\newcommand{\Msun}{M_{\odot}}
\newcommand{\Mstar}{M_{\star}}
\newcommand{\Mdyn}{M_{\mathrm{dyn}}}
\newcommand{\Mbar}{M_{\mathrm{bar}}}
\newcommand{\nbar}{\bar{n}}
\newcommand{\xiion}{\xi_{\mathrm{ion}}}
\newcommand{\RSVP}{\textsc{rsvp}}
\newcommand{\SIED}{\textsc{sied}}
\newcommand{\CLIO}{\textsc{clio}}

% ── Theorem-style environments ───────────────────────────────
\theoremstyle{definition}
\newtheorem{principle}{Principle}

% =============================================================
\begin{document}

% ── Title block ─────────────────────────────────────────────
\begin{center}
  {\LARGE\bfseries\sffamily Fossil Manifolds and Mediated Extraction}\\[0.55em]
  {\large\sffamily LAP1-B, Regulated Instability, and the Geometry of Observability}\\[1.2em]
  {\normalsize Flyxion}\\[0.4em]
  {\small\itshape Independent researcher}\\[0.3em]
  {\small May 2026}
\end{center}

\vspace{0.6em}

\epigraph{%
  The most informative structures in the universe are often both maximally
  preserved and maximally inaccessible.%
}{}

\vspace{0.8em}

% ── Abstract ────────────────────────────────────────────────
\begin{abstract}
\noindent
The recently reported JWST spectroscopy of LAP1-B --- an ultra-faint galaxy
at redshift $z \approx 6.6$ with oxygen abundance $(4.2 \pm 1.8)\times 10^{-3}$
solar, stellar mass below $3{,}300\,\Msun$, and dynamical mass of order
$10^{7}\,\Msun$ --- provides an unusual occasion for reflection on the
epistemology of preserved systems. This essay argues that LAP1-B exemplifies
a general principle: systems with the highest informational value are
frequently not the largest or most chemically evolved, but the least
overwritten. Beginning from observational facts, we develop this principle
through three interlocking operator classes. First, we characterize
\emph{emergence} through the band-limited instability formalism of Scalar
Irruption via Entropic Differential (\SIED{}), in which a preferred correlation
length $\ell \sim \sqrt{\kappa/|f''(\nbar)|}$ constrains which modes of
structure are physically admissible. Second, we analyze \emph{preservation}
through two distinct mechanisms: sparse internal history (the chemical event
log that remains readable because so few enrichment events have occurred) and
reionization-induced quenching (the temporal operator that froze the chemical
record by suppressing subsequent star formation). Third, we read the
gravitational lensing geometry as \emph{extraction}: the physical
amplification that makes an otherwise permanently inaccessible low-entropy
fossil detectable at cosmological distances. The essay concludes that the
deepest observational access to early-universe structure comes not through the
brightest or most massive objects, but through the rarest coincidence of
minimal overwriting and maximal geometric mediation --- a coincidence that
LAP1-B uniquely instantiates.
\end{abstract}

\vspace{1.0em}
\hrule
\vspace{1.4em}

% =============================================================
\section{Least Overwritten Systems}
\label{sec:overwrite}

Most of what we call chemical evolution in galaxies is more precisely chemical
averaging. Each successive generation of star formation stirs, dilutes, and
partially overwrites the abundance record left by previous generations.
Supernovae inject metals into a medium already enriched by earlier events;
inflows bring pristine gas that dilutes those metals; mergers mix populations
with distinct enrichment histories; turbulent transport spreads gradients that
would otherwise remain spatially coherent. The cumulative effect is that the
interstellar medium of a mature galaxy encodes not any single enrichment event
but an integral over many, and the integrand --- the specific sequence and
geometry of individual nucleosynthetic episodes --- is largely irrecoverable
from the integral alone.

This is not a merely practical limitation. It reflects something structural
about how information is stored in matter. A system that has undergone many
enrichment events has high chemical complexity but low historical resolution:
the state space accessible to it is large, but the trajectory that produced
its current state cannot be reconstructed without independent records of each
step. By contrast, a system that has undergone very few enrichment events
retains something closer to a first-order record of those events, precisely
because subsequent processing has not yet averaged the signal away. Its
chemical state is simpler, but its historical legibility is correspondingly
higher.

We propose calling such systems \emph{least overwritten} in the sense that
their current state retains the highest ratio of recoverable event history to
total elapsed time. Antiquity alone does not make a system least overwritten.
A very old galaxy with a rich star-formation history is thoroughly overwritten
despite its age. A young galaxy that formed early but evolved slowly ---
whether through low mass, environmental isolation, or reionization-induced
quenching --- may preserve far more of its initial conditions than a coeval
system that underwent vigorous enrichment. Least overwritten is therefore an
informational category, not a chronological one.

\begin{principle}[Informational Privilege of Low-Averaging Systems]
Among physically comparable systems, those that have undergone the fewest
averaging operations over their history retain the highest density of
recoverable information per observable degree of freedom.
\end{principle}

The corollary is that the instruments capable of reading such systems must
match the sparsity of the record: a detector calibrated for high-flux,
chemically complex sources will systematically miss the faintest, most
primitive signals. This creates a structural bias in observational surveys
toward overwritten systems and away from the very objects whose information
content is highest. We return to this selection problem in Section~\ref{sec:mediation}.

% =============================================================
\section{Event-History Legibility}
\label{sec:legibility}

A galaxy's abundance pattern is, in a precise sense, a projection of its
enrichment history onto the space of measurable ratios. The projection is
lossy: many distinct enrichment histories can produce similar abundance
patterns, and the projection operator --- the combined action of nucleosynthesis,
mixing, fallback, and chemical evolution --- discards information at every
step. Nevertheless, for a system that has undergone very few enrichment
events, the projection is substantially less lossy than for a chemically
mature system, because fewer intermediate states have intervened between the
first events and the currently observable state.

In sufficiently primitive systems, individual enrichment events may remain
legible not merely as integrated signatures but as geometrically specific
imprints. The carbon-to-oxygen ratio is a particularly sensitive indicator
because its value depends strongly on which fraction of each supernova's
yield survives to enter the subsequent gas phase. In energetic core-collapse
supernovae, both carbon and oxygen are produced and largely expelled into the
surrounding medium. In the so-called ``faint'' Population~III supernovae,
where the explosion energy is low enough that the oxygen-rich inner layers
undergo gravitational fallback into the central remnant while the carbon-rich
outer layers are ejected, the result is a characteristically elevated C/O
ratio at very low total metallicity \citep{Iwamoto2005, Nomoto2006}.

This is not merely a compositional signature. It is a record of a specific
dynamical asymmetry in the propagation of nucleosynthetic products: certain
outputs of stellar nucleosynthesis were prevented from entering the subsequent
evolutionary pathways of the system, not by random chance but by the physical
geometry of the collapse. The resulting abundance pattern therefore encodes
not only what elements were present but which physical channels remained open
for their propagation. In a chemically mature galaxy, repeated mixing would
erase much of this asymmetry. In a system that has experienced only one or a
few such events, the asymmetry survives as a nearly unprocessed record of
the earliest selective propagation.

It is in this sense that chemical abundances in primitive systems function as
an \emph{event log in matter}: a sparse, partially ordered record of
irreversible transitions whose outputs became the initial conditions for
everything that followed. The log is readable precisely because so little
has been written over it.

\begin{principle}[Event Logs in Matter]
In sufficiently low-averaging systems, observable abundance structures may
retain partially recoverable records of specific irreversible physical
transitions, allowing matter distributions to function as sparse historical
encodings of prior dynamical events.
\end{principle}

% =============================================================
\section{LAP1-B as a Preserved Enrichment Manifold}
\label{sec:lap1b}

\subsection{Chemical Primitivity and Low Averaging}

The JWST spectroscopy of LAP1-B \citep{Nakajima2026} presents a system
whose properties align with the least-overwritten category in a particularly
striking way. The galaxy is observed at redshift $z_{\mathrm{spec}} =
6.625 \pm 0.001$, corresponding to a cosmic age of approximately 800~Myr
after the Big Bang. It is essential to note that this redshift encodes an
observational epoch, not a contemporaneous state: the spectrum recorded by
JWST represents the condition of LAP1-B at the moment of photon emission,
not the galaxy's present cosmological state. The galaxy itself has almost
certainly undergone further evolution during the intervening propagation time.
What is preserved observationally is not the object as it currently exists,
but a temporally displaced projection carried by light through thirteen
billion years of cosmological propagation. Astronomy is, in this respect,
always the study of delayed observability.

The gas-phase oxygen abundance of LAP1-B is $(4.2 \pm 1.8) \times
10^{-3}$ solar --- well below the previously observed metallicity floor of
roughly $2\%$ solar --- and its elevated C/O ratio, at approximately one to
two times the solar value, is inconsistent with standard chemical enrichment
pathways but consistent with nucleosynthetic yields from faint Population~III
supernovae operating through the fallback mechanism described in
Section~\ref{sec:legibility}. The ionizing radiation field offers an
independent convergent constraint. The ionizing-photon production efficiency
satisfies $\log[\xiion] > 26.1$, a value incompatible with metal-enriched
stellar populations or black-hole accretion but consistent with either
metal-free Population~III stars or extreme Population~II stars with a
top-heavy initial mass function. These three independent observational lines
--- oxygen abundance, C/O ratio, and ionization hardness --- each point toward
the same conclusion: the stellar population responsible for LAP1-B's observable
state at the emission epoch is among the least chemically processed known.
Repeated enrichment events progressively transform discrete historical
signatures into integrated chemical averages whose originating trajectories
become increasingly irrecoverable; at the emission epoch, LAP1-B shows every
indication that this averaging process has barely begun.

\subsection{Dynamical Primitivity and Sparse Tracing}

The mass budget adds a structurally distinct dimension of primitivity.
The stellar mass is constrained by the absence of a detectable stellar
continuum to below $\Mstar < 3{,}300\,\Msun$ --- a quantity more typical
of a massive star cluster than a conventional galaxy --- while the dynamical
mass, inferred from H$\alpha$ emission-line broadening, is of order
$\Mdyn \sim 10^{7}\,\Msun$. The implied baryonic fraction falls below
approximately $1\%$. The luminous component is therefore not the dynamically
dominant constituent of the system; it is a sparse visible tracer of a much
larger organizing structure inferred only through its gravitational effects.

This dynamical primitivity is distinct from the chemical primitivity
discussed above, though the two reinforce each other. Chemical primitivity
reflects the absence of overwriting by successive nucleosynthetic episodes.
Dynamical primitivity reflects the absence of baryonic feedback: a stellar
mass of $\Mstar < 3{,}300\,\Msun$ cannot generate the supernovae, stellar
winds, or AGN activity needed to substantially reshape the dark matter
potential well it inhabits. The baryonic matter has not participated in
forming the gravitational structure it traces; it merely marks its minimum.
These two forms of primitivity together make LAP1-B an unusually clean
readout of both early enrichment geometry and early gravitational organization,
each preserved by the same underlying sparsity of its evolutionary history.

% =============================================================
\section{Constraint Basins and Sparse Tracers}
\label{sec:constraint}

\subsection{The Standard Halo Interpretation}

The standard $\Lambda$CDM interpretation of LAP1-B's mass budget treats the
dark matter halo as a pre-existing gravitational scaffold into which baryons
have condensed. The luminous stellar component is then understood as baryonic
matter that has settled to the potential minimum of this scaffold, with the
large dynamical-to-stellar mass ratio simply reflecting the efficiency with
which the halo assembled mass relative to the efficiency with which the
baryons converted gas into stars. On this view, the extreme mass ratio is
surprising in degree but not in kind: it represents the lower tail of a
continuous distribution of halo star-formation efficiencies, with LAP1-B
occupying the least efficient extreme.

\subsection{The Constraint-Basin Reinterpretation}

The following reinterpretation does not modify the underlying gravitational
dynamics. Rather, it changes which structural features are treated as
epistemologically primary.

An alternative descriptive framing, more natural within the \RSVP{}/\SIED{}
interpretive framework, inverts the phenomenological emphasis. Rather than
treating the dark matter halo as a container that the baryons fall into, one
can read the luminous matter as marking the deepest observable point of a
constraint basin --- by which we mean the dynamically admissible region toward
which trajectories converge under the effective organization of the system ---
whose full extent is inferred only dynamically. The baryons do not fill the
basin; they indicate where it is. They function as tracer particles in the
sense of fluid dynamics: sparse probes whose positions reveal the geometry of
a larger flow field that would otherwise be invisible.

Within the standard hierarchical picture, the dark matter halo functions
conceptually as a pre-existing container into which baryons subsequently
condense. The constraint-basin framing instead treats the luminous baryons
as sparse tracers marking the minima of an already-organized dynamical field.
$\Lambda$CDM takes the particle inventory as fundamental and derives the
organizing structure from it; the constraint-basin framing takes the
organizing structure as fundamental and treats the particle inventory as its
sparse observable projection.

\subsection{Epistemological Consequences of Sparse Tracing}

At LAP1-B's scale, the distinction between these framings is nearly
tractable in principle, precisely because the baryonic component is so sparse
that it genuinely cannot significantly reshape the potential well it inhabits.
In a massive galaxy, baryonic feedback --- stellar winds, supernova-driven
outflows, AGN jets --- substantially modifies the dark matter distribution on
timescales comparable to or shorter than the Hubble time, entangling the
constraint geometry and the particle content in ways that are difficult to
disentangle. In a system with $\Mstar < 3{,}300\,\Msun$, no such modification
has occurred. The baryons are tracers in a nearly pure sense: they probe the
organizing structure without participating in its formation.

The epistemological consequence is significant. In any system where the
visible matter has substantially modified the organizing structure, reading
the constraint geometry from the visible distribution requires a model of
the feedback history, introducing degeneracies that limit what can be inferred
about the primitive state. LAP1-B largely avoids this complication. The
visible matter is too sparse to have meaningfully altered what it traces.
This makes LAP1-B one of the cleanest available readouts of gravitational
organization in a state prior to substantial baryonic modification, and
therefore one of the few systems where the constraint-basin framing and the
standard halo framing differ in what they imply about the structure's
epistemic accessibility rather than merely about its description.

% =============================================================
\section{Regulated Instability and Preferred Irruption Scales}
\label{sec:sied}

\subsection{Hierarchical Accumulation and the Standard Picture}

The \SIED{} formalism offers a more precise account of why an object like
LAP1-B should exist at all and why it should exist at the scale it does.
In the standard hierarchical picture of structure formation, small-scale
structure forms first and subsequently merges into larger halos. The existence
of an isolated, extremely low-mass object at $z \approx 6.6$ is not, within
this picture, surprising --- it simply represents a halo at the low end of the
mass function. What is somewhat surprising is the object's apparent isolation
in chemical parameter space: its metallicity is not merely low but
discontinuously low relative to contemporaneous systems, occupying a regime
below what was previously identified as a metallicity floor.

\subsection{Threshold Instability and Spinodal Analogy}

The \SIED{} approach begins from a different structural assumption. Rather than
treating structure formation as hierarchical accumulation driven by gravitational
instability in an expanding medium, it treats it as a threshold instability in
a conserved scalar field. The \SIED{} instability mechanism is structurally
analogous to spinodal decomposition in condensed matter systems, where a
formerly homogeneous background loses local stability and spontaneously
separates into coherent domains once the curvature of the free-energy landscape
changes sign. The relevant criterion is therefore a sign change in the
free-energy curvature of the background state,
\begin{equation}
  f''(\nbar) < 0,
  \label{eq:spinodal}
\end{equation}
where $\nbar$ is the mean field density and $f$ is the local free-energy
functional. When condition~\eqref{eq:spinodal} is satisfied, the homogeneous
background becomes locally unstable to perturbations, and structure irrupts
into the formerly smooth medium --- not through gradual accumulation but
through a sudden local phase separation driven by the loss of convexity in
the free-energy landscape.

\subsection{Structural Rather Than Predictive Mathematics}

The mathematical relations introduced in this subsection are not presented
as quantitative fits to the specific observational parameters of LAP1-B.
Their function is instead structural and classificatory: they define a formal
vocabulary for regulated instability, preferred mode selection, and
finite-band irruption within a conserved medium. Readers should interpret
the equations that follow as an archetypal instability relation rather than
a direct predictive calculator for the observed dimensions of LAP1-B.

The crucial feature of \SIED{} that distinguishes it from unconstrained
gravitational collapse is the regulator term. The instability growth rate
takes the form
\begin{equation}
  \gamma(k) = a k^{2} - b k^{4},
  \label{eq:growthrate}
\end{equation}
where $k$ is the perturbation wavenumber and $a, b > 0$. Equation~\eqref{eq:growthrate}
should be understood as defining the structural logic of band-limited irruption,
not as a fitted model. The $ak^{2}$ term drives instability at all scales,
while the $bk^{4}$ term suppresses growth at high wavenumbers, preventing
arbitrary ultraviolet runaway. The regulated growth rate has a maximum at
$k_{\star} = \sqrt{a/2b}$, defining a preferred irruption scale. Below
$k_{\star}$, growth is too slow; above $k_{\star}$, growth is suppressed by
the regulator. Only a finite band of modes around $k_{\star}$ can irrupt
efficiently into visible structure.

The preferred spatial scale corresponding to this band is
\begin{equation}
  \ell \sim \sqrt{\frac{\kappa}{|f''(\nbar)|}},
  \label{eq:corrlen}
\end{equation}
where $\kappa$ is the gradient coefficient in the free-energy functional.
At present, the variables entering equation~\eqref{eq:corrlen} are not
empirically constrained by the available JWST dataset. The correspondence
developed here is qualitative and structural rather than numerically fitted;
the parameters $\kappa$ and $f''(\nbar)$ would need to be constrained against
independent observational data before any quantitative claim about
scale-selection could be made. This correlation length is not a free parameter
in principle, however: it is set by the physical properties of the background
state at the moment of irruption. Structure does not emerge at arbitrary scales
and subsequently filter down or up through a merger hierarchy; instead, it
emerges at a physically preferred scale determined by the local curvature of
the free-energy landscape.

\subsection{Interpretive Consequences for LAP1-B}

LAP1-B's unusual compactness and apparent isolation in chemical parameter
space are naturally interpreted within this framework as evidence of
constrained mode selection rather than incomplete hierarchical growth. The
object is not surprisingly small for its epoch; it is near the preferred
irruption scale for a localized instability in a medium still close to the
convexity boundary of its free-energy landscape. Its extreme chemical
primitivity corresponds to a region of the plenum that crossed the instability
threshold early and locally, forming a coherent low-complexity structure
before the surrounding medium had undergone substantial entropic processing.
In this reading, LAP1-B does not represent a fragment of a larger hierarchy
that failed to grow; it represents a preferred mode that irrupted near its
natural correlation length and has since been preserved against further
averaging by the very sparsity of its subsequent evolution.

\SIED{} does not provide an alternative microphysical account that would
observationally distinguish itself from $\Lambda$CDM at the level of LAP1-B's
reported data. What it offers is a different higher-level description of the
same phenomenology: a language in which scale-selection, constraint-basin
structure, and instability thresholds are primitive rather than derived.
The comparison with the observational findings is interpretive rather than
evidential. The argument is structural, not numerical.

% =============================================================
\section{Historical Encoding and the Admissibility of Propagation}
\label{sec:admissibility}

\subsection{Astrophysical Mechanism}

The fallback mechanism that produces LAP1-B's elevated C/O ratio is, at the
level of plasma physics, a straightforward dynamical process: the explosion
energy is insufficient to overcome the gravitational binding energy of the
oxygen-rich inner layers, which therefore fall back onto the central remnant
rather than being expelled into the surrounding medium. The carbon-rich outer
layers, less tightly bound and structurally positioned to receive the
initial momentum of the shock, escape into the interstellar gas. The result
is a characteristic chemical signature: elevated C/O at very low total
metallicity. This is the complete physical account, and it requires no
supplement from any non-standard framework.

\subsection{Interpretive Mapping onto Spherepop}

The preceding mechanism is entirely standard astrophysical plasma dynamics.
The Spherepop framework does not replace this account, but instead provides
a higher-order descriptive vocabulary for characterizing the selective
propagation structure implicit within it.

Mapping this fallback process onto the Spherepop framework, it instantiates
what can be called a \emph{propagation filter}: a physical constraint that
selectively allows certain products of an irreversible event to persist into
the state space available to the system's future evolution, while preventing
others from doing so. The oxygen-rich inner layers are dynamically refused
entry into the ambient gas; the carbon-rich outer layers are admitted. The
system's subsequent chemistry reflects not the full nucleosynthetic yield of
the event but the filtered yield --- those products that survived the
propagation constraint.

Spherepop does not treat this filter as operating through any semantic or
intentional mechanism. The refusal of oxygen to propagate is not a ``choice''
made by the system or by any abstract operator; it is the consequence of
specific initial conditions, explosion energetics, and gravitational geometry.
What Spherepop contributes is a vocabulary for describing the structural logic
of such constraints: the observation that physical systems routinely generate
irreversible events whose outputs are selectively admitted or blocked from
future evolutionary pathways, and that the pattern of admissions and blockings
constitutes a kind of historical record encoded in the system's current state.

In the case of LAP1-B, this record is unusually readable because so few
subsequent events have been written over it. The galaxy's current C/O ratio
is not merely a chemical datum; it is a partially preserved record of which
physical channels were open and which were closed during the earliest
nucleosynthetic episodes of its history. In a chemically mature galaxy,
successive rounds of enrichment would progressively overwrite this record,
replacing it with an integrated average over many such episodes. Here, the
record from one or at most a few events remains nearly intact. The galaxy's
chemistry functions as an event log, and the log is still sparse enough to
be read entry by entry rather than only as an average.

% =============================================================
\section{Spatial and Temporal Mediation}
\label{sec:mediation}

Gravitational lensing and cosmological reionization operate on primitive
systems in fundamentally different ways. Lensing is a spatial extraction
mechanism: it amplifies otherwise inaccessible signals across geometric
distance. Reionization, by contrast, is a temporal preservation mechanism:
it suppresses the subsequent evolutionary processes that would otherwise
overwrite primitive chemical records. Both are mediation operators in the
sense that each brings into the recoverable range a signal that would
otherwise remain permanently inaccessible, but they act on different
properties of the system and at different moments in its history.

\subsection{Gravitational Lensing as Spatial Extraction}

LAP1-B is observable at all only because the line of sight from the observer
to the galaxy passes through the gravitational lensing cluster
MACS~J0416.1-2403, which amplifies the galaxy's apparent flux by a factor
of approximately 100. Without this amplification, the galaxy would fall below
the detection threshold of any existing instrument, including JWST. The
lensing geometry is not a detail of the observation; it is the condition of
its possibility.

This circumstance has a straightforward observational explanation:
gravitational lensing by a massive foreground cluster bends and focuses
light from background sources, increasing their apparent brightness in
proportion to the solid angle subtended by the Einstein ring of the lens.
The amplification factor of $\sim 100$ corresponds to an effective increase
of five magnitudes, bringing a source that would be undetectable in unlensed
conditions into the range of spectroscopic accessibility. The physics is
well understood, and the lens model can be used to reconstruct the intrinsic
properties of the source from its lensed appearance.

What deserves emphasis is the structural role this plays in the epistemology
of observational access to primitive systems. The selection problem raised in
Section~\ref{sec:overwrite} --- that surveys systematically favor overwritten
systems because primitive systems are too faint to detect --- cannot be
resolved by improved detector sensitivity alone, at least not within the
current technological horizon. The intrinsic luminosity of a system like
LAP1-B scales with its stellar mass, and $\Mstar < 3{,}300\,\Msun$ produces
a surface brightness below the detection threshold of any foreseeable
direct-imaging survey at cosmological distances. The only currently available
mechanism for accessing such objects is geometric: the rare coincidence of a
suitable lensing mass along the line of sight.

The gravitational lens therefore functions, in this context, as a physical
sparse-inference operator: a mechanism that extracts a coherent signal from
a source that would otherwise be permanently inaccessible, by performing a
geometric amplification that brings the signal above the observational
threshold. The lens does not create information; it amplifies access to
information that exists but cannot be directly recovered. The operation is
analogous, in the language of \CLIO{} inference, to a compression-and-amplification
step that allows a low-entropy signal to be decoded from an attenuated projection.

\subsection{Reionization as Temporal Preservation}

A second mediation operator acts on LAP1-B's history rather than on its
visibility. The epoch of reionization --- the period during which ionizing
radiation from the first massive galaxies and quasars swept through the
neutral hydrogen of the early universe, heating and ionizing the intergalactic
medium --- acted as a quenching mechanism for low-mass halos. For halos with
deep gravitational wells, the heating was insufficient to drive gas escape,
and star formation continued. For halos as shallow as LAP1-B's, the intense
ultraviolet radiation boiled the cold, star-forming gas out of the potential
well or heated it above the threshold for continued condensation. Star
formation was permanently suppressed.

The same quenching processes that preserve primitive chemical structure by
suppressing subsequent star formation also reduce the luminosity of such
systems, rendering them progressively less accessible to direct observation.
This is the deeper inversion at the heart of the essay's argument:
reionization acted as an archive and as a blackout simultaneously. By
preventing further enrichment, it froze the chemical event log at a very
early entry. By suppressing star formation, it guaranteed that the object
emitting the frozen log would be too faint to observe without geometric
assistance. The fossil record is preserved by the same mechanism that renders
it inaccessible.

The fossil galaxies that survived reionization are precisely those that were
too small to sustain star formation against the ionizing radiation field, and
they are therefore also the least overwritten. Reionization paradoxically
preserves the fossil record by suppressing the enrichment that would overwrite
it. The epoch that is usually treated as a destructive and sterilizing
episode in galactic history is, for the smallest halos, a temporal
preservation operator: it drops the pin on the event log and ensures that
nothing further is written into it.

\subsection{The Coincidence of Three Independent Conditions}

The observation of LAP1-B thus depends on the simultaneous satisfaction of
three independent conditions: the galaxy must have remained chemically
primitive (low overwriting, governed by sparse early enrichment); it must
have remained dynamically tiny (low mass, enabling reionization quenching
and preventing baryonic feedback); and the line of sight must have passed
through a massive lensing cluster (geometric amplification across spatial
distance). The failure of any one condition would have made the observation
impossible. The galaxy's chemical primitivity and dynamical smallness are
not incidental properties that happen to make it interesting; they are the
same properties that make it nearly impossible to observe. The most
informationally rich objects are precisely the hardest to detect.

This inversion has consequences for how we understand the census of
early-universe structure. If LAP1-B is representative of a class rather
than an anomaly, the population of chemically primitive micro-galaxies in
the reionization era may be substantially larger than flux-limited surveys
suggest. Most of this population would be permanently inaccessible without
lensing assistance, not because the objects do not exist but because the
geometric coincidence required to observe them is rare. The fossil record of
early galaxy formation is likely more populated than our instruments can
directly see.

% =============================================================
\section{The Geometry of Observability}
\label{sec:clio}

\begin{principle}[Inverse Accessibility of Primitive Structure]
The systems retaining the highest density of recoverable historical
information are frequently those least accessible to direct observation,
because the sparsity that preserves their internal legibility also suppresses
their observational visibility.
\end{principle}

The preceding sections have assembled the components of a broader argument
about the structure of observational access to the early universe. We can
now state this argument in a unified form organized around three operator
classes.

\subsection{Emergence Through Regulated Instability}

Primitive systems --- those that are least overwritten in the sense of
Section~\ref{sec:overwrite} --- carry the highest density of recoverable
information about the earliest stages of structure formation. Their chemical
patterns encode event histories that more evolved systems have averaged away.
Their dynamical structures reflect organizing geometries that baryonic
feedback has not yet substantially modified. Their instability signatures,
where interpretable through frameworks like \SIED{}, reveal the preferred
scales at which structure first irrupts from a near-uniform background,
emerging not through gradual hierarchical accumulation but through threshold
phase separation at a physically mandated correlation length. In all these
respects, primitive systems are epistemically privileged relative to
chemically and dynamically mature ones.

\subsection{Preservation Through Quenching and Low Averaging}

And yet primitive systems are also maximally difficult to observe. Their
informational richness derives from the same sparsity that makes them
intrinsically faint. A galaxy with $\Mstar < 3{,}300\,\Msun$ is not merely
dim; it is, at cosmological distances, effectively invisible without geometric
assistance. The flux attenuation that comes from extreme distance and low
stellar mass conspires to place the most informative objects precisely below
the observational threshold. Reionization compounds this: by quenching star
formation and thereby preserving primitive chemical structure, it
simultaneously guarantees the low luminosity that makes such objects
undetectable without mediation.

\subsection{Extraction Through Geometric Mediation}

The consequence is that cosmological knowledge of the primitive epoch is
structurally dependent on mediated extraction. Direct access is not available.
What is available is access through operators --- physical, geometric, and
instrumental --- that amplify attenuated signals into the detectable range.
The \CLIO{} framework, which treats inference as the reconstruction of sparse
signals from attenuated projections, provides a natural language for
describing this structure. Observational cosmology in the primitive regime
is not a matter of gathering larger datasets from brighter sources; it is a
matter of identifying the rare configurations in which a mediation operator
--- a gravitational lens, a reionization geometry, an instrumental coincidence
--- brings a sparse, primitive signal into the recoverable range.

LAP1-B is not merely an early galaxy; it is a demonstration that such
configurations exist, and a proof of concept that the fossil record of the
earliest enrichment events can, under the right geometric conditions, be read
directly from the matter that preserves them.

The deepest observational access to the early universe, in this view, comes
not from the largest structures, which are bright but overwritten, nor from
the smallest structures in general, which are primitive but inaccessible,
but from the rarest coincidence: systems that are minimally overwritten,
sitting at preferred irruption scales within a constrained instability
spectrum, located at the intersection of a sufficient geometric mediation,
and preserved by the same quenching that froze their histories in place.
That coincidence is what LAP1-B instantiates. It is small enough to be a
fossil, primitive enough for its history to remain legible, and lensed
enough to be seen.

The broader implication extends beyond cosmology. Whenever historical
overwriting scales faster than observational accessibility, the systems
retaining the deepest information about primitive structure will also tend
to be the most difficult to observe directly. The problem of recovering
origins therefore becomes inseparable from the geometry of mediation itself:
the earliest histories survive not in the loudest systems, but in the
quietest ones whose signals remain barely recoverable across the intervening
layers of propagation, attenuation, and averaging.

% =============================================================
\section*{Acknowledgements}

This essay was written in response to the JWST spectroscopy reported by
\citet{Nakajima2026}. The interpretive frameworks developed here ---
\RSVP{}, \SIED{}, Spherepop, and \CLIO{} --- are theoretical constructs
under ongoing development and are not presented as empirically established
alternatives to standard cosmological models. The engagement with LAP1-B
is interpretive rather than evidential: the Nature paper's findings are
fully compatible with $\Lambda$CDM and are not claimed here as support for
any competing cosmological framework.

% =============================================================
\begin{thebibliography}{99}

\bibitem[Atek et~al.(2023)]{Atek2023}
Atek, H., Ellis, R., Richard, J., et~al.
\newblock Search for ultra-faint galaxies in the reionization era.
\newblock \textit{Monthly Notices of the Royal Astronomical Society}, 519(1):120--138, 2023.

\bibitem[Barkana \& Loeb(2001)]{Barkana2001}
Barkana, R., and Loeb, A.
\newblock In the beginning: the first sources of light and the reionization of the universe.
\newblock \textit{Physics Reports}, 349(2):125--238, 2001.

\bibitem[Bland-Hawthorn \& Freeman(2015)]{BlandHawthorn2015}
Bland-Hawthorn, J., and Freeman, K.
\newblock Near-field cosmology with dwarf elliptical galaxies.
\newblock \textit{Annual Review of Astronomy and Astrophysics}, 53:631--688, 2015.

\bibitem[Bromm \& Larson(2004)]{Bromm2004}
Bromm, V., and Larson, R.
\newblock The first stars.
\newblock \textit{Annual Review of Astronomy and Astrophysics}, 42:79--118, 2004.

\bibitem[Bromm \& Yoshida(2011)]{Bromm2011}
Bromm, V., and Yoshida, N.
\newblock The first galaxies.
\newblock \textit{Annual Review of Astronomy and Astrophysics}, 49:373--407, 2011.

\bibitem[Bullock \& Boylan-Kolchin(2017)]{Bullock2017}
Bullock, J.~S., and Boylan-Kolchin, M.
\newblock Small-scale challenges to the $\Lambda$CDM paradigm.
\newblock \textit{Annual Review of Astronomy and Astrophysics}, 55:343--387, 2017.

\bibitem[Carath\'{e}odory(1909)]{Caratheodory1909}
Carath\'{e}odory, C.
\newblock Untersuchungen \"{u}ber die Grundlagen der Thermodynamik.
\newblock \textit{Mathematische Annalen}, 67:355--386, 1909.

\bibitem[Colin et~al.(2019)]{Colin2019}
Colin, J., Mohayaee, R., Rameez, M., and Sarkar, S.
\newblock Evidence for anisotropy of cosmic acceleration.
\newblock \textit{Astronomy and Astrophysics}, 631:L13, 2019.

\bibitem[Dekel \& Silk(1986)]{Dekel1986}
Dekel, A., and Silk, J.
\newblock The origin of dwarf galaxies, cold dark matter, and biased galaxy formation.
\newblock \textit{Astrophysical Journal}, 303:39--55, 1986.

\bibitem[Eisenstein et~al.(2005)]{Eisenstein2005}
Eisenstein, D.~J., Zehavi, I., Hogg, D.~W., et~al.
\newblock Detection of the baryon acoustic peak in the large-scale correlation function
  of SDSS luminous red galaxies.
\newblock \textit{Astrophysical Journal}, 633(2):560--574, 2005.

\bibitem[Fan et~al.(2006)]{Fan2006}
Fan, X., Carilli, C., and Keating, B.
\newblock Observational constraints on cosmic reionization.
\newblock \textit{Annual Review of Astronomy and Astrophysics}, 44:415--462, 2006.

\bibitem[Ferrara(2012)]{Ferrara2012}
Ferrara, A.
\newblock The first galaxies and the likely discovery of their fossils in the Local Group.
\newblock \textit{Monthly Notices of the Royal Astronomical Society}, 421(1):282--294, 2012.

\bibitem[Frebel \& Norris(2015)]{Frebel2015}
Frebel, A., and Norris, J.~E.
\newblock Near-field cosmology with extremely metal-poor stars.
\newblock \textit{Annual Review of Astronomy and Astrophysics}, 53:631--688, 2015.

\bibitem[Furlanetto et~al.(2006)]{Furlanetto2006}
Furlanetto, S.~R., Oh, S.~P., and Briggs, F.
\newblock Cosmology at low frequencies: the 21\,cm transition and the high-redshift universe.
\newblock \textit{Physics Reports}, 433(4--6):181--301, 2006.

\bibitem[Gnedin(2000)]{Gnedin2000}
Gnedin, N.~Y.
\newblock Effect of reionization on structure formation in the universe.
\newblock \textit{Astrophysical Journal}, 542(2):535--541, 2000.

\bibitem[Greif(2015)]{Greif2015}
Greif, T.~H.
\newblock The numerical frontier of the high-redshift universe.
\newblock \textit{Computational Astrophysics and Cosmology}, 2(3):1--33, 2015.

\bibitem[Iwamoto et~al.(2005)]{Iwamoto2005}
Iwamoto, N., Nomoto, K., Umeda, H., et~al.
\newblock The first chemical enrichment in the universe and the formation of
  hyper metal-poor stars.
\newblock \textit{Science}, 309(5733):451--453, 2005.

\bibitem[Kang et~al.(2020)]{Kang2020}
Kang, Y., Lee, Y.-W., Kim, Y.-L., et~al.
\newblock Evidence for luminosity evolution of Type~Ia supernovae from their local environments.
\newblock \textit{Astrophysical Journal}, 889(1):8, 2020.

\bibitem[Kycia(2022)]{Kycia2022}
Kycia, R.
\newblock Entropy and ordering relations in thermodynamics.
\newblock \textit{Entropy}, 24(3):401, 2022.

\bibitem[Lieb \& Yngvason(1999)]{Lieb1999}
Lieb, E.~H., and Yngvason, J.
\newblock The physics and mathematics of the second law of thermodynamics.
\newblock \textit{Physics Reports}, 310(1):1--96, 1999.

\bibitem[Madau \& Dickinson(2014)]{Madau2014}
Madau, P., and Dickinson, M.
\newblock Cosmic star-formation history.
\newblock \textit{Annual Review of Astronomy and Astrophysics}, 52:415--486, 2014.

\bibitem[Maoz et~al.(2014)]{Maoz2014}
Maoz, D., Mannucci, F., and Nelemans, G.
\newblock Observational clues to the progenitors of Type~Ia supernovae.
\newblock \textit{Annual Review of Astronomy and Astrophysics}, 52:107--170, 2014.

\bibitem[Mo et~al.(2010)]{Mo2010}
Mo, H., van den Bosch, F., and White, S.
\newblock \textit{Galaxy Formation and Evolution}.
\newblock Cambridge University Press, Cambridge, 2010.

\bibitem[Nakajima et~al.(2026)]{Nakajima2026}
Nakajima, K., Ouchi, M., Zhang, Y., et~al.
\newblock An ultra-faint, chemically primitive galaxy forming in the reionization era.
\newblock \textit{Nature}, 653:363--367, 2026.

\bibitem[Navarro et~al.(1997)]{Navarro1997}
Navarro, J.~F., Frenk, C.~S., and White, S.~D.~M.
\newblock A universal density profile from hierarchical clustering.
\newblock \textit{Astrophysical Journal}, 490(2):493--508, 1997.

\bibitem[Nielsen et~al.(2016)]{Nielsen2016}
Nielsen, J.~T., Guffanti, A., and Sarkar, S.
\newblock Marginal evidence for cosmic acceleration from Type~Ia supernovae.
\newblock \textit{Scientific Reports}, 6:35596, 2016.

\bibitem[Nomoto et~al.(2006)]{Nomoto2006}
Nomoto, K., Tominaga, N., Umeda, H., et~al.
\newblock Nucleosynthesis yields of core-collapse supernovae and hypernovae,
  and galactic chemical evolution.
\newblock \textit{Nuclear Physics A}, 777:424--458, 2006.

\bibitem[{Planck Collaboration}(2020)]{Planck2020}
{Planck Collaboration}.
\newblock Planck 2018 results. VI. Cosmological parameters.
\newblock \textit{Astronomy and Astrophysics}, 641:A6, 2020.

\bibitem[Sarkar et~al.(2008)]{Sarkar2008}
Sarkar, S., Garcia-Bellido, J., and Haugb\o{}lle, T.
\newblock Is the evidence for dark energy secure?
\newblock \textit{International Journal of Modern Physics D}, 17(14):2505--2516, 2008.

\bibitem[Springel et~al.(2005)]{Springel2005}
Springel, V., White, S.~D.~M., Jenkins, A., et~al.
\newblock Simulations of the formation, evolution and clustering of galaxies and quasars.
\newblock \textit{Nature}, 435:629--636, 2005.

\bibitem[Tegmark et~al.(2004)]{Tegmark2004}
Tegmark, M., Strauss, M., Blanton, M., et~al.
\newblock Cosmological parameters from SDSS and WMAP.
\newblock \textit{Physical Review D}, 69(10):103501, 2004.

\bibitem[Tolstoy et~al.(2009)]{Tolstoy2009}
Tolstoy, E., Hill, V., and Tosi, M.
\newblock Star-formation histories, abundances, and kinematics of dwarf galaxies
  in the Local Group.
\newblock \textit{Annual Review of Astronomy and Astrophysics}, 47:371--425, 2009.

\bibitem[Umeda \& Nomoto(2003)]{Umeda2003}
Umeda, H., and Nomoto, K.
\newblock First-generation black-hole-forming supernovae and the metal abundance pattern
  of the most iron-poor stars.
\newblock \textit{Nature}, 422:871--873, 2003.

\bibitem[Weinberg(1989)]{Weinberg1989}
Weinberg, S.
\newblock The cosmological constant problem.
\newblock \textit{Reviews of Modern Physics}, 61(1):1--23, 1989.

\bibitem[White \& Rees(1978)]{White1978}
White, S.~D.~M., and Rees, M.
\newblock Core condensation in heavy halos: a two-stage theory for galaxy formation
  and clustering.
\newblock \textit{Monthly Notices of the Royal Astronomical Society}, 183:341--358, 1978.

\bibitem[Wise et~al.(2012)]{Wise2012}
Wise, J.~H., Turk, M.~J., Abel, T., and O'Shea, B.~W.
\newblock The birth of a galaxy: primordial metal enrichment and stellar populations.
\newblock \textit{Astrophysical Journal}, 745(1):50, 2012.

\bibitem[Yoshida et~al.(2008)]{Yoshida2008}
Yoshida, N., Omukai, T., Hernquist, L., and Springel, V.
\newblock Formation of primordial stars in a $\Lambda$CDM universe.
\newblock \textit{Science}, 321(5897):669--671, 2008.

\end{thebibliography}
\end{document}
